\appendix

\chapter{Numerical Implementation of the MLE algorithm}
\label{appendix:mle}
The fit procedure described in Sec. \ref{subsec:mle} employs a set of seven discrete matrices to evaluate the charge fractions $\vec{f}(x, y)$ needed for the calculation of the negative log-likelihood in Eq. \ref{eq:nll_multivariate}. Performing the minimization on discrete matrices could lead to numerical problems, due to the possibility of having non-physical zero gradients in case of small steps of the minimizer. This can prevent the algorithm from finding the global minimum, leading to errors in the reconstructed position and energy.

To avoid these problems, a bilinear interpolation of the seven charge fraction matrices is performed on the $(x, y)$ coordinates, making the evaluation of the $\operatorname{NLL}(x, y, Q)$ continuous and smooth across the pixel. For each pixel matrix, the interpolation is performed as follows:

\begin{itemize}[label=-]
\item calculate the fractional coordinates $(w_x, w_y)$ of $(x, y)$ with respect to the associated bin center and get the grid coordinates $(i_x, i_y)$ of the nearest lower-left bin;
\item get the values $(v_{00}, v_{10}, v_{01}, v_{11})$ of the four bins with grid coordinates $(i_x, i_y)$, $(i_x + 1, i_y)$, $(i_x, i_y + 1)$, $(i_x + 1, i_y + 1)$;
\item weight each bin value with the distance from $(x, y)$ to get the interpolated value:
\begin{equation}
f_i (x, y) = v_{00} (1 - w_x) (1 - w_y) + v_{10} w_x (1 - w_y) + v_{01} (1 - w_x) w_y + v_{11} w_x w_y.
\end{equation}
\end{itemize}

To speed up the minimizer and improve its stability, the gradient of the $\operatorname{NLL}$ with respect to the fit parameters $(x, y, Q)$ can be computed analytically. Recalling Eq. \ref{eq:nll_multivariate}, the negative log-likelihood can be written as a function of the generic parameter vector $\vec{\theta} = (x, y, Q)$, up to an additive constant, as
\begin{equation}
\operatorname{NLL}(\vec{\theta}) = \frac{1}{2} \vec{r}^{\,T} V^{-1} \vec{r} + \frac{1}{2} \ln |V|,
\end{equation}
where $\vec{r} = \vec{q} - Q \vec{f}(x, y)$ is the vector of charge residuals and $V$ is the covariance matrix defined in Eq. \ref{eq:cov_matrix}. Both the residuals vector $\vec{r}$ and the covariance matrix $V$ depend on the parameter vector $\vec{\theta}$.

The derivative with respect to a generic component $\theta_k$ is obtained by differentiating both the quadratic form and the log-determinant term
\begin{equation}
\frac{\partial \operatorname{NLL}}{\partial \theta_k} = \vec{r}^{\,T} V^{-1} \frac{\partial \vec{r}}{\partial \theta_k} - \frac{1}{2} \vec{r}^{\,T} V^{-1} \frac{\partial V}{\partial \theta_k} V^{-1} \vec{r} + \frac{1}{2} \operatorname{Tr} \left( V^{-1} \frac{\partial V}{\partial \theta_k} \right),
\label{eq:generic_nll_derivative}
\end{equation}
where it has been used that $\frac{\partial V^{-1}}{\partial \theta_k} = - V^{-1} \frac{\partial V}{\partial \theta_k} V^{-1}$ and $\frac{\partial \ln |V|}{\partial \theta_k} = \operatorname{Tr} \left( V^{-1} \frac{\partial V}{\partial \theta_k} \right)$.

Since the spatial dependence of the $\operatorname{NLL}$ is contained inside the charge fraction vector $\vec{f}(x, y)$, the spatial derivatives with respect to the coordinates $x$ and $y$ can be evaluated using the chain rule
\begin{equation}
\frac{\partial \operatorname{NLL}}{\partial x} = \sum_{i=0}^6 \frac{\partial \operatorname{NLL}}{\partial f_i} \frac{\partial f_i}{\partial x}, \quad \quad \frac{\partial \operatorname{NLL}}{\partial y} = \sum_{i=0}^6 \frac{\partial \operatorname{NLL}}{\partial f_i} \frac{\partial f_i}{\partial y},
\label{eq:nll_chain_rule}
\end{equation}
where the terms $\frac{\partial \operatorname{NLL}}{\partial f_i}$ are obtained by applying Eq. \ref{eq:generic_nll_derivative} with respect to each component $f_i$.

To avoid non-physical zero gradients and numerical issues during the minimization, the spatial derivatives of the charge fractions $\frac{\partial f_i}{\partial x}$ and $\frac{\partial f_i}{\partial y}$ are obtained by differentiating the bilinear interpolation formula
\begin{equation}
\begin{split}
\frac{\partial f_i}{\partial x} &= \frac{1}{dx} \left[ (v_{10} - v_{00}) (1 - w_y) + (v_{11} - v_{01}) w_y \right], \\
\frac{\partial f_i}{\partial y} &= \frac{1}{dy} \left[ (v_{01} - v_{00}) (1 - w_x) + (v_{11} - v_{10}) w_x \right],
\end{split}
\end{equation}
where $dx = dy$ is the square bin size.

Finally, the derivative with respect to the total charge $Q$ is evaluated directly from Eq. \ref{eq:generic_nll_derivative}, considering that $Q$ only enters the residuals $\vec{r}$ and the covariance matrix $V$
\begin{equation}
\frac{\partial \operatorname{NLL}}{\partial Q} = - \vec{f}^{\,T} V^{-1} \vec{r} - \frac{1}{2} \vec{r}^{\,T} V^{-1} \frac{\partial V}{\partial Q} V^{-1} \vec{r} + \frac{1}{2} \operatorname{Tr} \left( V^{-1} \frac{\partial V}{\partial Q} \right).
\end{equation}
